\title{xLSTM: Extended Long Short-Term Memory}

\begin{abstract}
During the 1990s, the foundational mechanisms of Long Short-Term Memory (LSTM) networks—namely, the constant error carousel and gating functions—were established. These innovations enabled LSTMs to play a pivotal role in the advancement of deep learning, particularly as the earliest architectures underlying Large Language Models (LLMs). Nonetheless, with the emergence of the Transformer architecture, which harnesses highly parallelizable self-attention, a paradigm shift occurred as Transformers began to outperform LSTMs at large scales. This development prompts us to investigate the capabilities of LSTMs when expanded to billions of parameters and enhanced with contemporary techniques inspired by modern LLMs, while addressing their known drawbacks. Our first contribution is the proposal of exponential gating, accompanied by suitable normalization and stabilization strategies. Additionally, we redesign the LSTM memory module, leading to the following variants: (i) sLSTM, which features a scalar memory cell, scalar update mechanism, and a novel memory mixing operation; and (ii) mLSTM, which achieves full parallelism via a matrix-based memory structure with an associated covariance-based update rule. By incorporating these advanced LSTM variants within residual block frameworks, we construct xLSTM blocks that can be composed in a stacked residual manner to form full xLSTM architectures. These enhancements—including exponential gating and restructured memory—substantially elevate the effectiveness of xLSTM models, enabling them to compare favorably with leading-edge Transformer and State Space Models, both in terms of scaling and overall performance.\\
Code available at: \url{https://github.com/NX-AI/xlstm}
\end{abstract}

\section{Introduction}
\label{sec:intro}

The Long Short-Term Memory (LSTM) ideas~\citep{Hochreiter:91,Hochreiter:97nips,Hochreiter:97}, i.e., the constant error carousel and gating, were introduced to overcome the vanishing gradient problem of recurrent neural networks~\citep{Hochreiter:91,Hochreiter:00book}:

\begin{align}
\label{eq:lstm_idea}
  c_t \ = \ f_t \ c_{t-1} \ + \ 
          i_t \  z_t \ , \quad  
          h_t \ = \ o_t \ \psi({c_t}) \ . 
\end{align}

The constant error carousel mechanism involves an additive modification of the cell state $c_{t-1}$ (shown in green), achieved through the incorporation of cell inputs $z_t$ and regulated by sigmoid gate functions (depicted in blue). The process is governed by the input gate $i_t$ and forget gate $f_t$, which together modulate the cell state's update, whereas the output of the memory cell—that is, the hidden state $h_t$—is determined by the output gate $o_t$. Before producing the hidden state, the cell state undergoes normalization or compression via the function $\psi$, after which the output gating yields the hidden state.

LSTMs have found success across a wide array of fields \citep{Hochreiter:01,Hochreiter:07,Schmidhuber:15}, dominating text generation until the emergence of Transformers in 2017~\citep{Vaswani:17}. Their robust performance has been established through applications in multiple sequence-focused tasks such as text generation~\citep{Graves:13,Karpathy:15blog}, handwriting synthesis~\citep{Graves:13}, sequence-to-sequence translation~\citep{Sutskever:14nips}, evaluation of computer programs~\citep{Zaremba:14arxiv}, producing image descriptions~\citep{Karpathy:15,Hossain:19}, source code creation~\citep{Karpathy:15blog}, rainfall-runoff forecasting~\citep{Kratzert:18,Kratzert:19}, as well as in hydrological systems for flood alerting~\citep{Nearing:24}. Within reinforcement learning, LSTMs continue to achieve state-of-the-art results for sequence modeling, as seen in the AlphaStar architecture for StarCraft II~\citep{Vinyals:17short}, the OpenAI Five agent for Dota 2~\citep{OpenAI:19}, and in systems that control magnetic fields for nuclear fusion~\citep{Degrave:22short}. LSTMs are particularly strong in acquiring high-level representations by effectively capturing and retaining semantic content in their memory cells~\citep{Karpathy:15blog}, which is illustrated by the discovery of number and syntax neurons~\citep{Lakretz:19}, linguistic neurons~\citep{Bau:19}, and sentiment neurons~\citep{Radford:17arxiv}. Despite the development of newer models, LSTMs continue to serve essential roles in demanding tasks~\citep{Degrave:22short,Nearing:24} and have proven their durability and continued relevance.

Although LSTMs have achieved remarkable results, they suffer from three primary shortcomings: (i) Inability to reconsider previously made storage choices. This weakness is illustrated through the {\it Nearest Neighbor Search} task (further details in Appendix~\ref{sec:appNearestNeighborSearch}): Given a reference vector, the sequence must be traversed to locate the vector most similar to the reference, retrieving its corresponding value at the end. The left side of Figure~\ref{fig:lstmProblems} presents mean squared error outcomes for this scenario. The standard LSTM faces difficulties in updating a value once stored, even if a more similar vector appears later; in contrast, our proposed xLSTM addresses this challenge by employing exponential gating mechanisms. (ii) Restricted storage ability, as all information must be compacted into scalar-valued cell states. This issue is highlighted using the {\it Rare Token Prediction} task. As demonstrated in the right panel of Figure~\ref{fig:lstmProblems}, the perplexity for predicting tokens in Wikitext-103~\citep{Merity:17}, separated by token frequency, shows that LSTM’s limited capacity leads to poorer performance on infrequent tokens. Our new xLSTM, equipped with a matrix-based memory, effectively overcomes this deficit. (iii) Inherent sequential dependence that impedes parallelism, a consequence of hidden-to-hidden connections linking successive time steps and requiring strictly sequential computation.

The emergence of Transformers~\citep{Vaswani:17} in the domain of language modeling is a direct response to these LSTM constraints. A natural question arises: What level of language modeling performance could be attained by resolving these LSTM limitations and scaling their architecture up to match the scale of current Large Language Models?

\section{Extended Long Short-Term Memory}
\label{sec:xLSTM}

To address the shortcomings of LSTM, the Extended Long Short-Term Memory (xLSTM) framework proposes two key enhancements to the conceptual basis presented in Equation~\eqref{eq:lstm_idea}. These enhancements—specifically, exponential gating and the introduction of advanced memory architectures—expand the LSTM family with two novel variants: (i) sLSTM (as detailed in Section~\ref{sec:sLSTM}), which employs a scalar-based memory system with scalar updates and incorporates memory mixing, and (ii) mLSTM (explained in Section~\ref{sec:mLSTM}), which utilizes matrix-based memory and updates via a covariance (outer product) mechanism that supports complete parallelism. Both variants incorporate exponential gating to augment the original LSTM formulation. The mLSTM, in particular, removes the traditional memory mixing (i.e., recurrent hidden-to-hidden interactions) to facilitate parallel computation. Additionally, both sLSTM and mLSTM architectures can be generalized to handle multiple memory cells. In this arrangement, sLSTM introduces inter-cell memory mixing, and can also be organized with multiple heads, such that there is no mixing between heads—only among cells within each head. This multi-head extension, together with exponential gating, creates a novel approach to inter-cell memory interaction. For mLSTM, configurations with several heads and cells are functionally indistinguishable.

By embedding these LSTM variants into residual block structures, one obtains xLSTM blocks (refer to Section~\ref{sec:xLSTMblock}). Stacking these blocks within deep networks yields xLSTM architectures (see Section~\ref{sec:xLSTMarchitecure} for more details). An illustration of the xLSTM architecture and its components is presented in Figure~\ref{fig:xlstm_sketch}.

\subsection{Review of the Long Short-Term Memory}
\label{sec:orgLSTM}

The foundational concept behind LSTM~\citep{Hochreiter:91,Hochreiter:97nips,Hochreiter:97} established the scalar memory cell as its core mechanism for information retention and manipulation, providing a solution to the vanishing gradient problem~\citep{Hochreiter:91,Hochreiter:00book} via the constant error carousel (i.e., the mechanism for updating the cell state). Three gating units—namely, input, output, and forget gates—regulate the flow of information within the memory cell. The introduction of the forget gate is credited to \citet{Gers:00}. At each time step $t$, the state of the LSTM memory cell evolves according to the following update equations:

\begin{align}
c_t \ &= \  f_t \ c_{t-1} \ + \ i_t \ z_t &  & &\text{cell state} \\
h_t  \ &= \ o_t \ \tilde{h}_t \ , & \tilde{h}_t \ &= \ \psi \left( c_t\right)
  &\text{hidden state} \\
z_t \ &= \ \varphi \left( \tilde{z}_t \right) \ , 
  &\tilde{z}_t \ &=  \ \mathbf{w}^\top_{z} \ \mathbf{x}_t \ + \
  r_{z}  h_{t-1} \ + \  b_{z} \ \
  &\text{cell input} \\
i_t \ &= \ \sigma \left( \tilde{i}_t  \right) \ , 
  &\tilde{i}_t \ &= \ \mathbf{w}^\top_{i} \ \mathbf{x}_t \ + \
  r_{i}  \ h_{t-1} \ + \  b_{i} \ \
  &\text{input gate} \\
f_t \ &= \ \sigma \left(  \tilde{f}_t \right) \ , 
  &\tilde{f}_t \ &= \ \mathbf{w}^\top_{f} \ \mathbf{x}_t  \ + \
  r_{f}  \ h_{t-1} \ + \  b_{f} \ \
  &\text{forget gate} \\
o_t \ &= \ \sigma \left( \tilde{o}_t \right) \ , 
  &\tilde{o}_t  \ &= \ \mathbf{w}^\top_{o} \ \mathbf{x}_t \ + \
  r_{o}  \ h_{t-1} \ + \  b_{o} \ \
  &\text{output gate} 
\end{align}

The vectors $\mathbf{w}_{z}$, $\mathbf{w}_{i}$, $\mathbf{w}_{f}$, and $\mathbf{w}_{o}$ denote the input weights connecting the input $\mathbf{x}_t$ to the cell input, input gate, forget gate, and output gate, respectively. The parameters $r_{z}$, $r_{i}$, $r_{f}$, and $r_{o}$ are the recurrent weights that link the previous hidden state $h_{t-1}$ to the cell input, as well as the respective gates. Associated bias terms are given by $b_{z}$, $b_{i}$, $b_{f}$, and $b_{o}$. The functions $\varphi$ and $\psi$ represent the activation mechanisms for the cell input and hidden state, with $\tanh$ being the standard choice in most cases. The activation function $\psi$ serves to constrain the values of the cell state, which otherwise could become unbounded. Sigmoid activations, that is, $\sigma \left( x \right)= 1/(1+\exp(-x))$, are applied to all gating units. Subsequent designs have aggregated several scalar memory cells into a vector form, $\mathbf{c}_t \in \mathbb{R}^{d}$, which permits the application of recurrent weight matrices $\mathbf{R} \in \mathbb{R}^{d \times d}$ for the purpose of integrating outputs across memory cells~\citep{Greff:15}; see further information in Appendix~\ref{sec:apporgLSTM}. Ablation experiments have indicated that each part of the memory cell plays a necessary role~\citep{Greff:15}.

\subsection{sLSTM}
\label{sec:sLSTM}

In order to provide LSTMs with the capacity to adjust their storage choices, we propose the use of exponential gates (indicated in red), supplemented by normalization and stabilization mechanisms. Specifically, the input and forget gates are permitted to employ exponential activation functions. For the purpose of normalization, we incorporate a normalizer state, which accumulates the sum obtained by multiplying the input gate by all subsequent forget gates.

The scalar sLSTM forward pass is:

\begin{align}
\label{eq:slstmforward}
c_t \ &= \  f_t \ c_{t-1} \ + \ i_t \ z_t & & 
 &\text{cell state} \\
n_t \ &= \  f_t \ n_{t-1} \ + \ i_t  & & 
 &\text{normalizer state} \\
h_t \ &= \ o_t \ \tilde{h}_t \ ,
  &\tilde{h}_t \ &= \ c_t / n_t
&\text{hidden state} \\
z_t \ &= \ \varphi \left( \tilde{z}_t \right) \ , 
  &\tilde{z}_t \ &=  \ \mathbf{w}^\top_{z} \ \mathbf{x}_t \ + \
  r_{z}  h_{t-1} \ + \  b_{z}
  &\text{cell input} \\
i_t \ &= \ \!\exp \left( \tilde{i}_t  \right) \ , 
  &\tilde{i}_t \ &= \ \mathbf{w}^\top_{i} \ \mathbf{x}_t \ + \
  r_{i}  \ h_{t-1} \ + \  b_{i}
  &\text{input gate} \\
f_t \ &= \ \sigma \left(  \tilde{f}_t \right) \ \text{OR} \
\!\exp \left(  \tilde{f}_t \right) \ ,
  &\tilde{f}_t \ &= \ \mathbf{w}^\top_{f} \ \mathbf{x}_t  \ + \
  r_{f}  \ h_{t-1} \ + \  b_{f}
  &\text{forget gate} \\
o_t \ &= \ \sigma \left( \tilde{o}_t \right) \ , 
  &\tilde{o}_t  \ &= \ \mathbf{w}^\top_{o} \ \mathbf{x}_t \ + \
  r_{o}  \ h_{t-1} \ + \  b_{o}
  &\text{output gate} 
\end{align}

We transfer the original LSTM gating techniques, i.e., input- and/or hidden-dependent gating plus bias term, to the new architectures. Exponential activation functions can lead to large values that cause overflows. Therefore, we stabilize gates with an additional state \mbox{$m_t$~\citep{Milakov:18arxiv}}:

\begin{align}
\label{eq:slstmstabil}
m_t \ &= \ \max \left( \log ( f_t ) + m_{t-1} , \log ( i_t ) \right) &\text{stabilizer state} \\
i'_t \ &= \ \!\exp \left( \log \left ( i_t \right) - m_t \right) \ = \ \!\exp \left( \tilde{i}_t - m_t \right) \ \, 
  &\text{stabil. input gate} \\
f'_t \ &= \ \!\exp \left( \log \left( f_t \right) + m_{t-1} - m_t \right) \ \, 
  &\text{stabil. forget gate}
\end{align}

As demonstrated in Appendix~\ref{sec:appsLSTM}, substituting $f_t$ with $f'_t$ and $i_t$ with $i'_t$ during the forward computation does not affect the network’s overall output nor the gradients of the loss with respect to its parameters.

\paragraph{New Memory Mixing.} Similar to the standard LSTM (refer to Appendix~\ref{sec:appsLSTM}), sLSTM is capable of utilizing several memory cells. The inclusion of multiple cells facilitates memory mixing through recurrent connections, denoted as $\mathbf{R}_{\mathbf{z}}$, $\mathbf{R}_{\mathbf{i}}$, $\mathbf{R}_{\mathbf{f}}$, and $\mathbf{R}_{\mathbf{o}}$, which connect the hidden state vector $\mathbf{h}$ to the memory cell input $\mathbf{z}$ as well as to the gating mechanisms $\mathbf{i}$, $\mathbf{f}$, and $\mathbf{o}$. A distinctive feature in this updated memory mixing process is the influence of exponential gating. The proposed sLSTM architecture allows for the existence of multiple heads, each supporting intra-head memory mixing, while no memory is shared across different heads. By introducing this notion of heads along with exponential gating, sLSTM presents a novel paradigm for mixing memory.

\subsection{mLSTM}
\label{sec:mLSTM}

In order to improve the memory capacity of LSTMs, we replace the standard scalar memory cell $c \in \mathbb{R}$ with a matrix memory cell $\mathbf{C} \in \mathbb{R}^{d \times d}$. As a result, accessing stored information now involves matrix multiplication. At each time step $t$, the model stores a key-value pair, where the key is denoted by $\mathbf{k}_t \in \mathbb{R}^d$ and the value by $\mathbf{v}_t \in \mathbb{R}^d$, adopting the nomenclature from Transformer architectures. At a later time point $t + \tau$, a query vector $\mathbf{q}_{t+\tau} \in \mathbb{R}^d$ is used to retrieve the corresponding value $\mathbf{v}_t$. This formulation follows the framework of Bidirectional Associative Memories (BAMs) as described in \citep{Kohonen:72,Anderson:72,Nakano:72,Anderson:77}.

The covariance update rule~\citep{Sejnowski:77,Dayan:91} for storing a key-value pair is

\begin{align}
\mathbf{C}_t \ &= \ \mathbf{C}_{t-1} \ + \ \mathbf{v}_{t} \ \mathbf{k}_t^\top \ .
\end{align}

We posit that applying a layer-norm to inputs prior to their projection onto keys and values ensures these vectors have zero mean. The covariance update rule achieves optimality~\citep{Dayan:91} in the sense of maximizing the separability among retrieved binary vectors, which is tantamount to optimizing the signal-to-noise ratio. If one restricts retrieval to pairwise interactions—thus adopting the quadratic computational cost characteristic of attention mechanisms~\citep{Krotov:16,Krotov:17,Ramsauer:21}—even greater separability can be attained. The covariance update rule aligns with the mechanism employed by Fast Weight Programmers~\citep{Schmidhuber:92ncfastweights,Schlag:21}, where subsequent modifications incorporated a constant decay rate (applied to $\mathbf{C}_{t-1}$) and a constant learning rate (applied to $\mathbf{v}_{t} \mathbf{k}_t ^\top$)~\citep{Ba:16}. Following this perspective, we embed the covariance update rule into the LSTM architecture, such that the forget gate serves as a decay rate, the input gate represents the learning rate, and the output gate modulates the magnitude of the retrieved vector.

For this matrix memory, the normalizer state is the weighted sum of key vectors, where each key vector is weighted by the input gate and all future forget gates. Again, the normalizer state keeps record of the strength of the gates. Since the dot product between query and normalizer state can be close to zero, we use the absolute value of this dot product and lower bound it by a threshold (typically 1.0) as done previously~\citep{Sun:23arxiv}. The mLSTM forward pass is:

\begin{align}
\label{eq:mlstm_recurrent_begin}
\mathbf{C}_t \ &= \  f_t \ \mathbf{C}_{t-1} \ + \ 
  i_t \ \mathbf{v}_t \ \mathbf{k}_t^\top & &  &\text{cell state} \\
\mathbf{n}_t \ &= \  f_t \ \mathbf{n}_{t-1} \ + \ 
  i_t \ \mathbf{k}_t & &  &\text{normalizer state} \\
\mathbf{h}_t  \ &= \ \mathbf{o}_t \ \odot \ \tilde{\mathbf{h}}_t
\ , \qquad \qquad 
\tilde{\mathbf{h}}_t \ = \   \mathbf{C}_t \mathbf{q}_t \ / \ 
  \max \left\{ |\mathbf{n}_t^\top \mathbf{q}_t|, 1 \right\} 
   &&&\text{hidden state} \\
\mathbf{q}_t \ &= \ \mathbf{W}_q \ \mathbf{x}_t \ + \ \mathbf{b}_q  & &  &\text{query input} \\
\mathbf{k}_t \ &= \ \frac{1}{\sqrt{d}} \mathbf{W}_k \ \mathbf{x}_t \ + \ \mathbf{b}_k  & &  &\text{key input} \\
\mathbf{v}_t \ &= \ \mathbf{W}_v \ \mathbf{x}_t \ + \ \mathbf{b}_v  & &  &\text{value input} \\
i_t \ &= \ \!\exp \!  \! \left( \tilde{i}_t  \right) \ , 
 \qquad \qquad \ \ \ \ \,
  \tilde{i}_t \ = \ \mathbf{w}^\top_{i} \ \mathbf{x}_t \ + \  b_{i} \ \ \ 
  &&&\text{input gate} \\
f_t \ &= \ \sigma \!  \left(  \tilde{f}_t \right) \ \text{OR} \ \!\exp \!  \! \left(  \tilde{f}_t \right) , \ \ \: \!
\tilde{f}_t \ = \ \mathbf{w}^\top_{f} \ \mathbf{x}_t  \ + \
  b_{f}
  &&& \text{forget gate} \\
\label{eq:mlstm_recurrent_end}
\mathbf{o}_t \ &= \ \sigma \left( \tilde{\mathbf{o}}_t \right) \ , \qquad \qquad \quad \ \ \ \,
  \tilde{\mathbf{o}}_t  \ = \ \mathbf{W}_{\mathbf{o}} \ \mathbf{x}_t \ + \
  \mathbf{b}_{\mathbf{o}}
  &&&\text{output gate} 
\end{align}

The architecture of mLSTM allows for several memory cells, similar to the standard LSTM. In mLSTM, employing multiple heads is functionally equivalent to using multiple memory cells due to the absence of memory interaction across cells. To maintain the stability of mLSTM's exponential gating mechanisms, we apply the same stabilization methods described for sLSTM (refer to Equation~\ref{eq:slstmstabil}). Because mLSTM does not involve the intermingling of memory, it is possible to express its recurrence relations in a parallelized manner. Additional information can be found in Appendix~\ref{sec:appmLSTM}.

\subsection{xLSTM Architecture}
\label{sec:xLSTMarch}

\paragraph{xLSTM Blocks.}
\label{sec:xLSTMblock}
An xLSTM block is designed to capture past information through non-linear transformations in a high-dimensional feature space, thereby enabling more effective distinction between diverse historical sequences or contexts. Such distinction is fundamental for accurately forecasting subsequent elements in a sequence, such as predicting the next token. This approach is guided by Cover’s Theorem~\citep{Cover:65}, which asserts that non-linearly projected data are more likely to be linearly separable in higher-dimensional spaces than in their original, lower-dimensional form. We investigate two variants of residual block structures: (i) A residual block with a post up-projection (similar to the approach employed in Transformers), where the block initially performs a non-linear aggregation of the prior context in the input space, then projects the result into a higher-dimensional space via a linear operation, applies a non-linear activation, and subsequently linearly maps back into the original input space; refer to the left panel of Figure~\ref{fig:Backbones}
as well as the third column in Figure~\ref{fig:xlstm_sketch}. 
A comprehensive depiction can be found in Figure~\ref{fig:appsLSTM_detailed}
in the appendix. (ii) A residual block with a pre up-projection (analogous to architectures found in State Space Models), where an initial linear transformation projects to a higher-dimensional space,

The method performs a nonlinear aggregation of previous information within a high-dimensional context, followed by a linear transformation that projects the output back into the initial feature space. In scenarios where the xLSTM block incorporates an sLSTM, a post up-projection module is employed. Conversely, if the block contains an mLSTM, a pre up-projection module is chosen, as the memory size is increased in the high-dimensional setting. Further elaboration can be found in the left side of Figure~\ref{fig:Backbones}, the third panel of Figure~\ref{fig:xlstm_sketch}, or in Figure~\ref{fig:appsLSTM_detailed} within the appendix.

\paragraph{xLSTM Architecture. }
\label{sec:xLSTMarchitecure}
The xLSTM model is formed through the residual stacking of modular components, as outlined in~\citep{Srivastava:15,He:16}. This design follows the prevalent architecture of pre-LayerNorm~\citep{Ba:16b} residual backbones found in modern Large Language Models. Reference for the design can be found in the final panel of Figure~\ref{fig:xlstm_sketch}.

\subsection{Memory and Speed Considerations}

Unlike Transformers, xLSTM networks maintain linear computational cost and fixed memory requirements regardless of sequence length. Owing to their compressive memory, xLSTM architectures are highly compatible with real-world deployment and edge computing scenarios.

The mLSTM architecture utilizes a parameter-free memory structure, but it incurs significant computational overhead due to the use of a $d \times d$ matrix for both storage and updates. This design represents a compromise between enhancing memory capacity and managing computational demands. However, these matrix operations are inherently parallelizable on GPUs, which helps to minimize their impact on actual processing time.

Whereas mLSTM supports parallelization akin to techniques like FlashAttention~\citep{Dao:22,Dao:23} or GLA~\citep{Yang:23arxiv}, sLSTM is inherently sequential as a result of its memory mixing via hidden-to-hidden connections, which hinders parallel processing. Nonetheless, we engineered an efficient CUDA implementation for sLSTM with register-level GPU memory enhancements, resulting in a runtime that is typically at most twice that of mLSTM.

\section{Related Work}
\label{sec:related}

\paragraph{Linear Attention.} Numerous approaches have been proposed to address the Transformer’s quadratic time complexity relative to context length, with the goal of achieving attention mechanisms whose computation scales linearly. For instance, the Synthesizer bypasses token-to-token dependencies by learning synthetic attention weights~\citep{Tay:20arxiv}. Linformer approximates self-attention through a low-rank projection, resulting in a linearized computation~\citep{Wang:20arxiv}. The Linear Transformer transforms the attention operation itself into a linear function~\citep{Katharopoulos:20}. Performer employs a positive orthogonal random features technique to approximate the attention softmax function with linear complexity~\citep{Choromanski:21}. Alternatives to attention include substituting it with efficient long-range convolutional operators, as demonstrated in Structured Global Convolution (SGConv)~\citep{Li:22} and the Hyena Hierarchy~\citep{Poli:23}.

\paragraph{State Space Models.} State Space Models (SSMs) have surged in popularity recently, largely due to their linear complexity with respect to sequence length and strong empirical results in comparison to Transformer architectures. Among the initial models introduced in this line was the Structured State Space sequence model (S4)~\citep{Gu:21}. This was followed by several variants, including the Diagonal State Space (DSS) model~\citep{Gupta:22}, Gated State Space (GSS) models~\citep{Mehta:22}, S5 model~\citep{Smith:22}, Bidirectional Gated SSM (BiGS)~\citep{Wang:22}, H3 model~\citep{Fu:23}, and Mamba~\citep{Gu:24arxiv}.

\paragraph{Recurrent Neural Networks.} Recurrent Neural Networks (RNNs) have been proposed as alternatives to the Transformer architecture and attention mechanisms, largely because they offer linear computational scaling with respect to context length. Notably, RNNs that utilize Deep Linear Recurrent Units (LRUs) have yielded strong outcomes in language modeling tasks~\citep{Orvieto:23, De:24arxiv}, as have both the Hierarchically Gated Linear RNN (HGRN)~\citep{Qin:23} and its successor HGRN2~\citep{Qin:24arxiv}. Among approaches targeting large language models, the RWKV method~\citep{Peng:23arxivshort,Peng:24arxivshort} stands out, demonstrating performance on par with that of Transformer-based systems.

\paragraph{Gating.}
The concept of gating, foundational to LSTM architectures, has resurfaced and found new interpretations in several recent models. Notably, gating mechanisms are a core component in HGRN~\citep{Qin:23}, HGRN2~\citep{Qin:24arxiv}, Gated Linear Attention (GLA)~\citep{Yang:23arxiv}, Gated State Space (GSS) models~\citep{Mehta:22}, Bidirectional Gated SSM (BiGS)~\citep{Wang:22}, Moving Average Equipped Gated Attention (MEGA)~\citep{Ma:22}, as well as in RWKV~\citep{Peng:23arxivshort} and Mamba~\citep{Gu:24arxiv}.

\paragraph{Covariance Update Rule.} To bolster the storage capabilities of the mLSTM cell, we introduced a matrix-based memory utilizing a covariance update rule. This update strategy is also adopted in several alternative frameworks, including Fast Weight Programmers \citep{Schmidhuber:92ncfastweights,Schlag:21}, RWKV-5 and RWKV-6~\citep{Peng:24arxivshort}, Retention~\citep{Sun:23arxiv}, Linear Transformer~\citep{Katharopoulos:20}, and HGRN2~\citep{Qin:24arxiv}.

\paragraph{Most Related.} The models most conceptually aligned with xLSTM include Retention~\citep{Sun:23arxiv}, RWKV~\citep{Peng:23arxivshort,Peng:24arxivshort}, and HGRN2~\citep{Qin:24arxiv}. All these models incorporate mechanisms involving matrix memory or gating. Nonetheless, unlike the novel sLSTM, these methods lack support for memory mixing. The introduction of memory mixing is crucial for addressing state tracking challenges, making LSTMs more powerful than State Space Models (SSMs) and Transformers~\citep{Merrill:24,Deletang:23}. State tracking is essential in tasks such as program evaluation and entity tracking over long narrative sequences.

\paragraph{Residually Stacking Architectures.} The xLSTM architectures, in line with the majority of modern large-scale deep learning systems, are formulated through the residual stacking of foundational components~\citep{Srivastava:15,He:16}.

This framework is foundational not only for deep convolutional neural networks~\citep{He:16}, but also for the architecture of Transformers~\citep{Vaswani:17}. Transformers, in turn, are the essential mechanism underlying the development of Large Language Models (LLMs) such as GPT-3~\citep{Brown:20short}, ChatGPT~\citep{Schulman:22short}, GPT-4~\citep{Achiam:23gpt4short}, Megatron-LM~\citep{Shoeybi:19}, Gopher~\citep{Rae:21short}, ERNIE 3.0 Titan~\citep{Wang:21arxivshort}, GLaM~\citep{Du:21glamshort}, Chinese M6~\citep{Lin:21}, the multilingual AlexaTM 20B~\citep{Soltan:22}, OPT~\citep{Zhang:22opt}, Chinchilla~\citep{Hoffmann:22short}, BLOOM~\citep{Scao:22arxivshort}, GLM-130B~\citep{Zeng:22arxivshort}, LaMDA~\citep{Thoppilan:22short}, PaLM~\citep{Chowdhery:22short}, Llama~\citep{Touvron:23arxiv}, and Gemini~\citep{GeminiTeam:23arxiv,Reid:24arxiv}.

\section{Experiments}
\label{sec:Exp}

This section presents an empirical analysis of xLSTM, emphasizing its comparison with established techniques for language modeling. Section~\ref{sec:ExpSyntethic} examines the unique strengths of xLSTM through synthetic benchmarks. In Section~\ref{sec:ExpComparison}, we report validation set perplexities for multiple state-of-the-art language modeling approaches, each trained on 15B tokens from the SlimPajama corpus~\citep{Soboleva:23}. For xLSTM, we further perform a set of ablation experiments using this dataset. Subsequently, we analyze how each model scales, adopting evaluation methodologies similar to those of \citet{Kaplan:20} and \citet{Brown:20short}.

Section~\ref{sec:ExpLanguage} details an extensive language modeling study. Here, xLSTM and the top-performing models from Section~\ref{sec:ExpComparison} are trained on a larger dataset comprising 300B SlimPajama tokens~\citep{Soboleva:23}. The comparison proceeds along several axes: evaluating the models' capacity for generalization to longer sequences, examining their performance in terms of validation perplexity and on select downstream tasks~\citep{Sutawika:24}, benchmarking across 571 text domains from the PALOMA language suite~\citep{Magnusson:23arxivshort}, and revisiting the issue of scaling but now leveraging a dataset twenty times the size.

\label{sec:xlstm_blocks_notation}
Throughout our experiments, we denote xLSTM[$a$:$b$] to indicate that the composition of xLSTM blocks follows a ratio of $a/b$ between mLSTM-based and sLSTM-based blocks. As an illustration, xLSTM[7:1] specifies a configuration of eight total blocks, where seven are built upon mLSTM and a single block utilizes sLSTM. In the scenario where the total count of blocks is set to 48, this configuration corresponds to 42 mLSTM-based blocks and 6 sLSTM-based blocks. Additionally, all experimental setups consistently include pre up-projection blocks for mLSTM and post up-projection blocks for sLSTM.

\subsection{Synthetic Tasks and Long Range Arena}
\label{sec:ExpSyntethic}
We begin by evaluating the impact of xLSTM's exponential gating and memory mixing mechanisms using formal language tasks~\citep{Deletang:23}. Subsequently, we examine how well the newly introduced matrix memory in xLSTM performs on the Multi-Query Associative Recall benchmark~\citep{Arora:23arxiv}. Lastly, we analyze xLSTM’s capabilities for handling extended sequence data by measuring its results on the Long Range Arena benchmark~\citep{Tay:21}.

\paragraph{Evaluation of xLSTM’s Exponential Gating Incorporating Memory Mixing.} We assess the effectiveness of xLSTM’s novel exponential gating combined with memory mixing, a mechanism designed to address state tracking challenges~\citep{Merrill:24, Merrill:23}. The formal language benchmarks from \citet{Deletang:23} are adapted and broadened to include multi-length training, thereby supporting evaluation on length generalization. A comprehensive overview of all experimental tasks, as well as further analysis, is included in Appendix~\ref{sec:appExpSynthetic-formalLang}. Our evaluation contrasts xLSTM with competing architectures such as Transformers, State Space Models, and classic Recurrent Neural Networks. Accuracy is measured for only those tokens integral to each task, and results are normalized so that 0 reflects random guessing and 1 corresponds to perfect prediction. The specific configurations assessed include 2-block architectures: xLSTM[0:1] (pure sLSTM), xLSTM[1:0] (pure mLSTM), xLSTM[1:1], along with Llama, Mamba, RWKV, Retention, Hyena, LSTM, and LSTM incorporated in Transformer layers (LSTM (Block)). Outcomes from these assessments are summarized in Figure~\ref{fig:formal-main}. Architectures like standard Transformers or State Space Models that lack memory mixing—and, thus, effective state tracking—fail on tasks demanding such capabilities, such as recognizing regular grammars exemplified by the parity test. These observations support previous work asserting the fundamental limitations of Transformers and State Space models relative to RNNs~\citep{Merrill:24, Merrill:23, Deletang:23}.

\paragraph{Test of xLSTM's Memory Capacities on Associative Recall Tasks.} This study evaluates the memory storage abilities of xLSTM’s novel matrix-based memory using the Multi-Query Associative Recall task~\citep{Arora:23arxiv}. In each test sequence, a random assortment of key-value pairs is sampled from a sizable vocabulary, and the model must reliably store and later recall them. To further increase the difficulty compared to prior setups, we scale the number of key-value items as high as 256 and extend the sequence context to 2048 tokens. This design enables a more extensive assessment of model memory capabilities. We benchmark two-block configurations from Llama, Mamba, RWKV-5, RWKV-6, xLSTM[1:1], and xLSTM[1:0], judging their performance by their accuracy in retrieving key-value associations. With their exponentially large memory (relative to coding dimension)~\citep{Ramsauer:21}, Transformer models such as Llama serve as the reference for optimal performance. Figure~\ref{fig:mqar-main} presents comparative results. xLSTM[1:1] consistently outperforms all other non-Transformer candidates, including in small model settings. Notably, adding the sLSTM block enhances, rather than hampers, overall memory retention; this effect is especially notable for the hardest test involving 256 key-value pairs. Further experiments and extrapolation studies are described in Appendix~\ref{sec:appExpSynthetic-mqar}, suggesting that xLSTM's expanded memory capacities extend generalization to sequence contexts longer than those included in the training data.

\paragraph{Test of xLSTM's Long Context Capabilities on Long Range Arena.} In order to evaluate xLSTM's effectiveness with lengthy input sequences and expansive contexts, we conducted comparisons with various approaches on the Long Range Arena~\citep{Tay:21}. Across all tasks, xLSTM achieves robust and reliable results, indicating that its architecture is highly adept at managing diverse challenges associated with processing long-range contexts.

For more details, see Appendix~\ref{sec:appExpSynthetic-lra}.

\subsection{Method Comparison and Ablation Study}
\label{sec:ExpComparison}

In order to investigate the central question posed in this work—specifically, the capabilities of our newly proposed LSTM variants when applied to large-scale language modeling tasks—we conduct experiments by training xLSTMs, Transformers, State Space Models, as well as various alternative architectures, all utilizing 15B tokens sourced from SlimPajama under a consistent auto-regressive framework. We assess these models’ performances on the validation set and carry out ablation experiments focused on the xLSTMs.

\paragraph{Comparing xLSTM to Other Methods.} To enable a meaningful comparison, we trained all models using a dataset of 15B tokens drawn from SlimPajama~\citep{Soboleva:23}. Evaluation was carried out by measuring perplexity on the validation split. The following set of approaches are considered in our analysis: xLSTM (the novel approach proposed here), GPT-3 (Transformer)~\citep{Brown:20short}, Llama (Transformer)~\citep{Touvron:23arxiv}, H3 (SSM)~\citep{Fu:23}, Mamba (SSM)~\citep{Gu:24arxiv}, RWKV-4 (RNN)~\citep{Peng:23arxivshort}, RWKV-5 (RNN)~\citep{Peng:24arxivshort}, RWKV-6 (RNN)~\citep{Peng:24arxivshort}, GLA (linear Transformer)~\citep{Yang:23arxiv}, HGRN (RNN)~\citep{Qin:23}, HGRN2 (RNN)~\citep{Qin:24arxiv}, RetNet (linear Transformer)~\citep{Sun:23arxiv}, Hyena (linear Transformer)~\citep{Poli:23}, xLSTM[1:0], and xLSTM[7:1] (see Section~\ref{sec:xlstm_blocks_notation} for these notations). Training across all models used mixed precision; for RWKV-5, RWKV-6, GLA, and HGRN2, mixed precision was implemented without relying on PyTorch’s automated mixed precision (refer to Appendix Section~\ref{sec:appGenTrainProc}). 

We group the methods into (a) Transformers, (b) State Space Models (SSMs), and (c) Recurrent Neural Networks (RNNs), incorporating linear Transformers alongside the RNNs due to the linear attention mechanism replacing standard Transformer attention. Each model is configured to closely resemble a 350M-parameter GPT-3 model, employing an embedding size of 1024 and consisting of 24 residual blocks. Notably, only GPT-3 utilizes shared token and output embedding weights, resulting in a lower total parameter count. Table~\ref{tab:spaj15b_model_results} presents the validation perplexity scores, clearly demonstrating that xLSTM achieves superior performance relative to prior approaches. Further information is available in Appendix~\ref{sec:appExpComparison}. 

The scaling characteristics for these experiments are illustrated in Figure~\ref{fig:temp_sclaw15B}, and these results suggest that xLSTM retains its advantage at larger scales.

\paragraph{Ablation Studies.}
As shown in Table~\ref{tab:spaj15b_model_results} and Figure~\ref{fig:temp_sclaw15B}, xLSTM demonstrates superior performance in language modeling tasks when trained on 15B tokens from SlimPajama. To systematically analyze the contribution of each architectural modification from LSTM to xLSTM, we incrementally convert a standard LSTM model into the xLSTM configuration. The process begins by incorporating LSTM layers within pre-LayerNorm residual networks. Next, we adapt the architecture to utilize a post up-projection block. The final modification involves implementing exponential gating mechanisms and introducing matrix memory.

The results are shown in Table~\ref{tab:ablstudies} (top). 

Ablation experiments demonstrate that the marked gains in performance stem from a combination of exponential gating and the matrix memory. Furthermore, recognizing the crucial role of gating within RNNs and State Space Models, we examine various gating designs. As presented in Table~\ref{tab:ablstudies} (bottom), making each gate adaptive and dependent on the input yields progressively better results. Further analysis of specific backbone elements can be found in Appendix~\ref{sec:appAblDetails}.

\subsection{xLSTM as Large Language Model}
\label{sec:ExpLanguage}

This investigation concludes with extensive language modeling experiments aimed at evaluating the suitability of xLSTM as a large language model (LLM). To this end, we scale up the training by utilizing 300B tokens sourced from SlimPajama, aligning our dataset size with prior works such as Mamba~\citep{Gu:24arxiv} and Griffin~\citep{De:24arxiv}.

For empirical comparison, xLSTM is benchmarked against RWKV-4, Llama, and Mamba, each representing a distinct class of architectures outlined in Section~\ref{sec:ExpComparison}. RWKV-4 is chosen to exemplify recurrent neural networks, since suitable training precision configurations for RWKV-5, RWKV-6, and HGRN2 were determined only after the onset of the 300B-token training phase (see Appendix~\ref{sec:appGenTrainProc}). 

Models of various scales (125M, 350M, 760M, and 1.3B parameters) are trained and assessed for their ability to extrapolate to longer sequences and validated on a held-out dataset. Further evaluation covers downstream task performance, language modeling capability on the 571 text domains present in the PALOMA benchmark, and a thorough analysis of scaling law characteristics.

\paragraph{Sequence Length Extrapolation.}
\label{sec:spaj300B_ctx_extrapolation}
To begin, we assess the ability of large-scale (1.3B parameter) variants of xLSTM, RWKV-4, Llama, and Mamba to extrapolate to longer sequence lengths. Each model undergoes training using a fixed context window of 2048 tokens and is subsequently evaluated on contexts extending up to 16384 tokens. The comparative performance is depicted in Figure~\ref{fig:spaj300B_seq_extrapolation}. Notably, among the evaluated architectures, xLSTM exhibits consistently low perplexity even as the test sequence length increases.

\paragraph{Validation Perplexity and Downstream Tasks.}
\label{sec:spaj_downstream_eval}
Next, we investigate all evaluated model scales, measuring xLSTM, RWKV-4, Llama, and Mamba performance on the SlimPajama validation set for next-token prediction, as well as on various downstream benchmarks emphasizing common sense reasoning abilities.

The third column of Table~\ref{tab:lmeval} presents the perplexity scores on the validation set for each method. Among the evaluated models, both xLSTM[1:0] and xLSTM[7:1] consistently achieve the lowest perplexity across all considered model sizes. The remaining columns in Table~\ref{tab:lmeval} summarize results on various downstream benchmarks. Across nearly all tasks and model capacities, xLSTM outperforms alternative approaches—with the exception of certain instances in the ARC task, where Mamba demonstrates superior results. Comprehensive results and further explanation can be found in Appendix~\ref{sec:appExpLanguage}.

\paragraph{Performance on PALOMA Language Tasks.} Third, we evaluate the next token prediction accuracy of xLSTM, RWKV-4, Llama, and Mamba across all model sizes using the PALOMA language benchmarks~\citep{Magnusson:23arxivshort}. Specifically, we assess performance by calculating perplexity scores over 571 different text domains, encompassing sources from nytimes.com to the r/depression subreddit on Reddit. 

Table~\ref{tab:ppleval} provides perplexity values organized by broad language modeling (first seven columns) and by more specialized, domain-specific benchmarks (last five columns). On these language tasks, xLSTM[1:0] consistently yields better results than xLSTM[7:1]. 

Notably, xLSTM[1:0] achieves a lower perplexity compared to Mamba in 568 of the 571 text domains (99.5\%), outperforms Llama in 486 out of 571 domains (85.1\%), and surpasses RWKV-4 in 570 of the 571 domains (99.8\%). Further analysis is available in Appendix~\ref{sec:appExpLanguage}.

\paragraph{Scaling Laws.} Fourth, we investigate the power-law scaling trends, which provide insight into how model performance can be predicted for larger parameter sizes~\citep{Kaplan:20,Brown:20short}. Figure~\ref{fig:temp_sclaw300B} illustrates these scaling patterns. While all models exhibit comparable scaling curves, their performance baselines differ. RWKV-4 demonstrates the weakest results, followed in order by Llama and Mamba. In comparison, xLSTM outperforms Mamba by a gap similar to that between Mamba and Llama. These results suggest that xLSTM is likely to maintain, or even increase, its advantage in performance relative to Transformers and State-Space models as model size grows.

\textbf{Generation Times and Maximal Throughput.} We assess the time taken for text generation in Figure~\ref{fig:inference_time_generation} as well as the peak throughput rates in Figure~\ref{fig:inference_time_decoding_throughput} (left) for various xLSTM variants, all at the 1.3B parameter scale. Our analysis includes a comparison with similarly sized Mamba, Llama, and RWKV models implemented in HuggingFace, noting the use of a static key-value cache for the Llama architecture. At the point when these experiments were conducted, compiling the complete cache for the Transformer model or employing \texttt{torch.compile} on the Mamba implementation was unsuccessful. In our generation benchmarks, each model is evaluated at a batch size of 1 and a pre-fill length of 16 tokens—settings that particularly benefit the Transformer. 

The results depicted in Figure~\ref{fig:inference_time_generation} reveal that xLSTM, together with other recurrent approaches like Mamba and RWKV-4, exhibits linear scaling characteristics, whereas Llama demonstrates quadratic scaling behavior. To evaluate decoding throughput, we tested various batch sizes and prefill levels for the Llama model. As shown in Figure~\ref{fig:inference_time_decoding_throughput} (right), xLSTM takes advantage of its constant memory usage, allowing support for much larger batch sizes than Llama and consequently attaining superior throughput performance.

\section{Limitations}

(i) Unlike mLSTM, the sLSTM’s approach to memory mixing hinders the use of parallelizable computations, which in turn limits rapid parallel execution. Despite this, we implemented an efficient CUDA kernel for sLSTM, achieving runtimes that are under twice as slow compared to our parallel mLSTM version. (ii) The mLSTM CUDA kernels are currently unoptimized; as a result, our present implementation operates approximately four times slower than either FlashAttention or the scan-based mechanism employed in Mamba. Improved speed could be achieved with more optimized CUDA kernels, drawing inspiration from FlashAttention. (iii) The computational demands of the mLSTM’s matrix memory are significant since it processes $d \times d$ matrices. Nevertheless, both the updating and accessing of memory remain parameter-free and are amenable to parallelization with conventional matrix operations, so the practical increase in wall-clock computation time remains modest. (iv) Careful consideration is necessary when initializing the forget gates. (v) Because the matrix memory does not grow with sequence length, increasing sequence lengths can cause memory saturation with larger context windows. However, we have not observed this as a significant constraint for sequence lengths up to 16k, as discussed in Section~\ref{sec:spaj300B_ctx_extrapolation}. (vi) Owing to the resource-intensive nature of training large language models, neither the network design nor hyperparameter choices have been thoroughly fine-tuned—especially for larger xLSTM models. It is likely that comprehensive optimization is essential for xLSTM to realize its maximum effectiveness.

\section{Conclusion}
We have provided a partial response to our initial question: To what extent can LSTM-based language modeling be advanced by scaling up to billions of parameters? Our findings suggest that such scaling can at minimum match the capabilities of leading methods such as Transformers and State Space Models. By introducing exponential gating with memory mixing and devising a novel memory architecture, we have transformed the LSTM into the xLSTM. The xLSTM demonstrates strong comparative results in language modeling tasks, standing shoulder-to-shoulder with state-of-the-art approaches like Transformers and State Space Models. Scaling law analysis suggests that expanding the size of xLSTM models will position them as formidable rivals to current Large Language Models based on Transformer architectures. Furthermore, xLSTM exhibits promise for substantial contributions in a variety of other domains, including Reinforcement Learning, Time Series Prediction, and the modeling of physical systems.

\end{document}